\section{Criteria for flatness}
\label{section-criteria-flatness}

\noindent
In this section we prove some important technical lemmas in the Noetherian
case. We will (partially) generalize these to the non-Noetherian case
in Section \ref{section-more-flatness-criteria}.

\begin{lemma}
\label{lemma-mod-injective}
Suppose that $R \to S$ is a local homomorphism of local rings with $S$
Noetherian.
Denote $\mathfrak m$ the maximal ideal of $R$. Let $M$ be a flat $R$-module
and $N$ a finite $S$-module. Let $u : N \to M$ be a map of $R$-modules.
If $\overline{u} : N/\mathfrak m N \to M/\mathfrak m M$
is injective then $u$ is injective.
In this case $M/u(N)$ is flat over $R$.
\end{lemma}

\begin{proof}
First we claim that $u_n : N/{\mathfrak m}^nN \to M/{\mathfrak m}^nM$
is injective for all $n \geq 1$. We proceed by induction, the base
case is that $\overline{u} = u_1$ is injective. By our assumption that $M$
is flat over $R$ we have  a short exact sequence
$0 \to M \otimes_R {\mathfrak m}^n/{\mathfrak m}^{n + 1}
\to M/{\mathfrak m}^{n + 1}M \to M/{\mathfrak m}^n M \to 0$.
Also, $M \otimes_R {\mathfrak m}^n/{\mathfrak m}^{n + 1}
= M/{\mathfrak m}M \otimes_{R/{\mathfrak m}}
{\mathfrak m}^n/{\mathfrak m}^{n + 1}$. We have
a similar exact sequence $N \otimes_R {\mathfrak m}^n/{\mathfrak m}^{n + 1}
\to N/{\mathfrak m}^{n + 1}N \to N/{\mathfrak m}^n N \to 0$
for $N$ except we do not have the zero on the left. We also
have $N \otimes_R {\mathfrak m}^n/{\mathfrak m}^{n + 1}
= N/{\mathfrak m}N \otimes_{R/{\mathfrak m}}
{\mathfrak m}^n/{\mathfrak m}^{n + 1}$. Thus the map $u_{n + 1}$ is
injective as both $u_n$ and the map
$\overline{u} \otimes \text{id}_{{\mathfrak m}^n/{\mathfrak m}^{n + 1}}$ are.

\medskip\noindent
By Krull's intersection theorem
(Lemma \ref{lemma-intersect-powers-ideal-module-zero})
applied to $N$ over the ring $S$ and the ideal $\mathfrak mS$
we have $\bigcap \mathfrak m^nN = 0$. Thus the injectivity
of $u_n$ for all $n$ implies $u$ is injective.

\medskip\noindent
To show that $M/u(N)$ is flat over $R$, it suffices to show that
$\text{Tor}_1^R(M/u(N), R/I) = 0$ for every ideal $I \subset R$,
see Lemma \ref{lemma-characterize-flat}. From the short exact sequence
$$
0 \to N \xrightarrow{u} M \to M/u(N) \to 0
$$
and the flatness of $M$ we obtain an exact sequence of Tors
$$
0 \to \text{Tor}_1^R(M/u(N), R/I) \to N/IN \to M/IM
$$
See Lemma \ref{lemma-long-exact-sequence-tor}. Thus it suffices to show that
$N/IN$ injects into $M/IM$. Note that $R/I \to S/IS$ is a local
homomorphism of local rings with $S/IS$ Noetherian, $N/IN \to M/IM$ is a map
of $R/I$-modules, $N/IN$ is finite over $S/IS$, and $M/IM$ is flat over
$R/I$ and $u \bmod I : N/IN \to M/IM$ is injective modulo $\mathfrak m$.
Thus we may apply the first part of the proof to $u \bmod I$ and we conclude.
\end{proof}

\begin{lemma}
\label{lemma-grothendieck}
Suppose that $R \to S$ is a flat and local ring homomorphism of Noetherian
local rings. Denote $\mathfrak m$ the maximal ideal of $R$.
Suppose $f \in S$ is a nonzerodivisor in $S/{\mathfrak m}S$.
Then $S/fS$ is flat over $R$, and $f$ is a nonzerodivisor in $S$.
\end{lemma}

\begin{proof}
Follows directly from Lemma \ref{lemma-mod-injective}.
\end{proof}

\begin{lemma}
\label{lemma-grothendieck-regular-sequence}
Suppose that $R \to S$ is a flat and local ring homomorphism of Noetherian
local rings. Denote $\mathfrak m$ the maximal ideal of $R$.
Suppose $f_1, \ldots, f_c$ is a sequence of elements of
$S$ such that the images $\overline{f}_1, \ldots, \overline{f}_c$
form a regular sequence in $S/{\mathfrak m}S$.
Then $f_1, \ldots, f_c$ is a regular sequence in $S$ and each
of the quotients $S/(f_1, \ldots, f_i)$ is flat over $R$.
\end{lemma}

\begin{proof}
Induction and Lemma \ref{lemma-grothendieck}.
\end{proof}

\begin{lemma}
\label{lemma-free-fibre-flat-free}
Let $R \to S$ be a local homomorphism of Noetherian
local rings. Let $\mathfrak m$ be the maximal
ideal of $R$. Let $M$ be a nonzero finite $S$-module.
Suppose that (a) $M/\mathfrak mM$
is a free $S/\mathfrak mS$-module, and (b) $M$ is flat over $R$.
Then $M$ is free and $S$ is flat over $R$.
\end{lemma}

\begin{proof}
Let $\overline{x}_1, \ldots, \overline{x}_n$ be a basis
for the free module $M/\mathfrak mM$. Choose
$x_1, \ldots, x_n \in M$ with $x_i$ mapping to $\overline{x}_i$. Let
$u : S^{\oplus n} \to M$ be the map which maps the $i$th
standard basis vector to $x_i$. By Lemma \ref{lemma-mod-injective}
we see that $u$ is injective. On the other hand, by
Nakayama's Lemma \ref{lemma-NAK} the map is surjective. The
lemma follows.
\end{proof}

\begin{lemma}
\label{lemma-complex-exact-mod}
Let $R \to S$ be a local homomorphism of local Noetherian
rings. Let $\mathfrak m$ be the maximal ideal of $R$.
Let $0 \to F_e \to F_{e-1} \to \ldots \to F_0$
be a finite complex of finite $S$-modules. Assume that
each $F_i$ is $R$-flat, and that the complex
$0 \to F_e/\mathfrak m F_e \to F_{e-1}/\mathfrak m F_{e-1}
\to \ldots \to F_0 / \mathfrak m F_0$ is exact.
Then $0 \to F_e \to F_{e-1} \to \ldots \to F_0$
is exact, and moreover the module
$\Coker(F_1 \to F_0)$ is $R$-flat.
\end{lemma}

\begin{proof}
By induction on $e$. If $e = 1$, then this is exactly
Lemma \ref{lemma-mod-injective}. If $e > 1$, we see
by Lemma \ref{lemma-mod-injective} that $F_e \to F_{e-1}$
is injective and that $C = \Coker(F_e \to F_{e-1})$
is a finite $S$-module flat over $R$. Hence we can
apply the induction hypothesis to the complex
$0 \to C \to F_{e-2} \to \ldots \to F_0$.
We deduce that $C \to F_{e-2}$ is injective
and the exactness of the complex follows, as well
as the flatness of the cokernel of $F_1 \to F_0$.
\end{proof}

\noindent
In the rest of this section we prove two versions of what is called the
``{\it local criterion of flatness}''. Note also the interesting
Lemma \ref{lemma-CM-over-regular-flat} below.

\begin{lemma}
\label{lemma-prepare-local-criterion-flatness}
Let $R$ be a local ring with maximal ideal $\mathfrak m$
and residue field $\kappa = R/\mathfrak m$.
Let $M$ be an $R$-module. If $\text{Tor}_1^R(\kappa, M) = 0$,
then for every finite length $R$-module $N$ we have
$\text{Tor}_1^R(N, M) = 0$.
\end{lemma}

\begin{proof}
By descending induction on the length of $N$.
If the length of $N$ is $1$, then $N \cong \kappa$
and we are done. If the length of $N$ is more than
$1$, then we can fit $N$ into a short exact sequence
$0 \to N' \to N \to N'' \to 0$ where $N'$, $N''$ are
finite length $R$-modules of smaller length.
The vanishing of $\text{Tor}_1^R(N, M)$ follows
from the vanishing of $\text{Tor}_1^R(N', M)$
and $\text{Tor}_1^R(N'', M)$ (induction hypothesis)
and the long exact sequence of Tor groups, see Lemma
\ref{lemma-long-exact-sequence-tor}.
\end{proof}

\begin{lemma}[Local criterion for flatness]
\label{lemma-local-criterion-flatness}
Let $R \to S$ be a local homomorphism of local Noetherian
rings. Let $\mathfrak m$ be the maximal ideal of $R$,
and let $\kappa = R/\mathfrak m$.
Let $M$ be a finite $S$-module. If $\text{Tor}_1^R(\kappa, M) = 0$,
then $M$ is flat over $R$.
\end{lemma}

\begin{proof}
Let $I \subset R$ be an ideal. By Lemma \ref{lemma-flat} it suffices
to show that $I \otimes_R M \to M$ is injective. By Remark
\ref{remark-Tor-ring-mod-ideal} we see that this kernel is
equal to $\text{Tor}_1^R(M, R/I)$. By
Lemma \ref{lemma-prepare-local-criterion-flatness}
we see that $J \otimes_R M \to M$ is injective for all ideals
of finite colength.

\medskip\noindent
Choose $n >> 0$ and consider the following short exact
sequence
$$
0
\to I \cap \mathfrak m^n
\to I \oplus \mathfrak m^n
\to I + \mathfrak m^n
\to 0
$$
This is a sub sequence of the short exact sequence
$0 \to R \to R^{\oplus 2} \to R \to 0$. Thus we get the diagram
$$
\xymatrix{
(I\cap \mathfrak m^n) \otimes_R M \ar[r] \ar[d] &
I \otimes_R M \oplus \mathfrak m^n \otimes_R M \ar[r] \ar[d] &
(I + \mathfrak m^n) \otimes_R M \ar[d] \\
M \ar[r] &
M \oplus M \ar[r] &
M
}
$$
Note that $I + \mathfrak m^n$ and $\mathfrak m^n$
are ideals of finite colength.
Thus a diagram chase shows that
$\Ker((I \cap \mathfrak m^n)\otimes_R M \to M)
\to \Ker(I \otimes_R M \to M)$
is surjective. We conclude in particular that
$K = \Ker(I \otimes_R M \to M)$ is contained
in the image of $(I \cap \mathfrak m^n) \otimes_R M$
in $I \otimes_R M$. By Artin-Rees, Lemma \ref{lemma-Artin-Rees}
we see that $K$ is contained
in $\mathfrak m^{n-c}(I \otimes_R M)$ for some $c > 0$
and all $n >> 0$. Since $I \otimes_R M$ is a finite
$S$-module (!) and since $S$ is Noetherian, we see
that this implies $K = 0$. Namely, the above implies
$K$ maps to zero in the $\mathfrak mS$-adic completion
of $I \otimes_R M$. But the map from $S$
to its $\mathfrak mS$-adic completion is faithfully flat
by Lemma \ref{lemma-completion-faithfully-flat}.
Hence $K = 0$, as desired.
\end{proof}

\noindent
In the following we often encounter the conditions
``$M/IM$ is flat over $R/I$ and $\text{Tor}_1^R(R/I, M) = 0$''.
The following lemma gives some consequences of these conditions
(it is a generalization of
Lemma \ref{lemma-prepare-local-criterion-flatness}).

\begin{lemma}
\label{lemma-what-does-it-mean}
Let $R$ be a ring.
Let $I \subset R$ be an ideal.
Let $M$ be an $R$-module.
If $M/IM$ is flat over $R/I$ and $\text{Tor}_1^R(R/I, M) = 0$ then
\begin{enumerate}
\item $M/I^nM$ is flat over $R/I^n$ for all $n \geq 1$, and
\item for any module $N$ which is annihilated by $I^m$ for some $m \geq 0$
we have $\text{Tor}_1^R(N, M) = 0$.
\end{enumerate}
In particular, if $I$ is nilpotent, then $M$ is flat over $R$.
\end{lemma}

\begin{proof}
Assume $M/IM$ is flat over $R/I$ and $\text{Tor}_1^R(R/I, M) = 0$.
Let $N$ be an $R/I$-module. Choose a short exact sequence
$$
0 \to K \to \bigoplus\nolimits_{i \in I} R/I \to N \to 0
$$
By the long exact sequence of $\text{Tor}$ and the vanishing of
$\text{Tor}_1^R(R/I, M)$ we get
$$
0 \to \text{Tor}_1^R(N, M) \to K \otimes_R M \to
(\bigoplus\nolimits_{i \in I} R/I) \otimes_R M \to N \otimes_R M \to 0
$$
But since $K$, $\bigoplus_{i \in I} R/I$, and $N$ are all annihilated
by $I$ we see that
\begin{align*}
K \otimes_R M & = K \otimes_{R/I} M/IM, \\
(\bigoplus\nolimits_{i \in I} R/I) \otimes_R M & =
(\bigoplus\nolimits_{i \in I} R/I) \otimes_{R/I} M/IM, \\
N \otimes_R M & = N \otimes_{R/I} M/IM.
\end{align*}
As $M/IM$ is flat over $R/I$ we conclude that
$$
0 \to K \otimes_{R/I} M/IM \to
(\bigoplus\nolimits_{i \in I} R/I) \otimes_{R/I} M/IM \to
N \otimes_{R/} M/IM \to 0
$$
is exact. Combining this with the above we conclude that
$\text{Tor}_1^R(N, M) = 0$ for any $R$-module $N$ annihilated by $I$.

\medskip\noindent
Let us prove (2) by induction on $m$. The case $m = 1$ was
done in the previous paragraph. For $N$ annihilated by $I^m$
for $m > 1$ we may choose an exact sequence
$0 \to N' \to N \to N'' \to 0$ with $N'$ and $N''$ annihilated
by $I^{m - 1}$. For example one can take $N' = IN$ and $N'' = N/IN$.
Then the exact sequence
$$
\text{Tor}_1^R(N', M) \to
\text{Tor}_1^R(N, M) \to
\text{Tor}_1^R(N'', M)
$$
and induction prove the vanishing we want.

\medskip\noindent
Finally, we prove (1). Given $n \geq 1$ we have to show that $M/I^nM$
is flat over $R/I^n$. In other words, we have to show that the functor
$N \mapsto N \otimes_{R/I^n} M/I^nM$ is exact on the category of $R$-modules
$N$ annihilated by $I^n$. However, for such $N$ we have
$N \otimes_{R/I^n} M/I^nM = N \otimes_R M$. By the vanishing of
$\text{Tor}_1$ in (2) we see that the functor $N \mapsto N \otimes_R M$
is exact on the category of $N$ annihilated by some power of $I$
and we conclude.
\end{proof}

\begin{lemma}
\label{lemma-what-does-it-mean-again}
Let $R$ be a ring. Let $I \subset R$ be an ideal.
Let $M$ be an $R$-module.
\begin{enumerate}
\item If $M/IM$ is flat over $R/I$ and $M \otimes_R I/I^2 \to IM/I^2M$
is injective, then $M/I^2M$ is flat over $R/I^2$.
\item If $M/IM$ is flat over $R/I$ and $M \otimes_R I^n/I^{n + 1}
\to I^nM/I^{n + 1}M$ is injective for $n = 1, \ldots, k$,
then $M/I^{k + 1}M$ is flat over $R/I^{k + 1}$.
\end{enumerate}
\end{lemma}

\begin{proof}
The first statement is a consequence of
Lemma \ref{lemma-what-does-it-mean} applied with $R$ replaced by $R/I^2$
and $M$ replaced by $M/I^2M$ using that
$$
\text{Tor}_1^{R/I^2}(M/I^2M, R/I) =
\Ker(M \otimes_R I/I^2 \to IM/I^2M),
$$
see Remark \ref{remark-Tor-ring-mod-ideal}.
The second statement follows in the same manner using induction
on $n$ to show that $M/I^{n + 1}M$ is flat over $R/I^{n + 1}$ for
$n = 1, \ldots, k$. Here we use that
$$
\text{Tor}_1^{R/I^{n + 1}}(M/I^{n + 1}M, R/I^n) =
\Ker(M \otimes_R I^n/I^{n + 1} \to I^nM/I^{n + 1}M)
$$
for every $n$.
\end{proof}

\begin{lemma}[Variant of the local criterion]
\label{lemma-variant-local-criterion-flatness}
Let $R \to S$ be a local homomorphism of Noetherian
local rings. Let $I \not = R$ be an ideal in $R$.
Let $M$ be a finite $S$-module. If $\text{Tor}_1^R(M, R/I) = 0$
and $M/IM$ is flat over $R/I$, then $M$ is flat over $R$.
\end{lemma}

\begin{proof}
First proof: By
Lemma \ref{lemma-what-does-it-mean}
we see that $\text{Tor}_1^R(\kappa, M)$ is zero where $\kappa$
is the residue field of $R$. Hence we see that $M$
is flat over $R$ by
Lemma \ref{lemma-local-criterion-flatness}.

\medskip\noindent
Second proof: Let $\mathfrak m$ be the maximal ideal of $R$.
We will show that $\mathfrak m \otimes_R M \to M$ is injective,
and then apply
Lemma \ref{lemma-local-criterion-flatness}.
Suppose that $\sum f_i \otimes x_i \in \mathfrak m \otimes_R M$
and that $\sum f_i x_i = 0$ in $M$. By the equational criterion
for flatness Lemma \ref{lemma-flat-eq} applied to $M/IM$
over $R/I$ we see there exist $\overline{a}_{ij} \in R/I$
and $\overline{y}_j \in M/IM$ such that
$x_i \bmod IM = \sum_j \overline{a}_{ij} \overline{y}_j $
and $0 = \sum_i (f_i \bmod I) \overline{a}_{ij}$.
Let $a_{ij} \in R$ be a lift of $\overline{a}_{ij}$ and
similarly let $y_j \in M$ be a lift of $\overline{y}_j$.
Then we see that
\begin{eqnarray*}
\sum f_i \otimes x_i
& = &
\sum f_i \otimes x_i +
\sum f_ia_{ij} \otimes y_j -
\sum f_i \otimes a_{ij} y_j
\\
& = &
\sum f_i \otimes (x_i - \sum a_{ij} y_j) +
\sum (\sum f_i a_{ij}) \otimes y_j
\end{eqnarray*}
Since $x_i - \sum a_{ij} y_j \in IM$ and
$\sum f_i a_{ij} \in I$ we see that there exists
an element in $I \otimes_R M$ which maps to our given
element $\sum f_i \otimes x_i$ in $\mathfrak m \otimes_R M$.
But $I \otimes_R M \to M$ is injective by assumption (see
Remark \ref{remark-Tor-ring-mod-ideal}) and we win.
\end{proof}

\noindent
In particular, in the situation of
Lemma \ref{lemma-variant-local-criterion-flatness}, suppose that
$I = (x)$ is generated by a single element $x$ which is
a nonzerodivisor in $R$. Then $\text{Tor}_1^R(M, R/(x)) = (0)$
if and only if $x$ is a nonzerodivisor on $M$.

\begin{lemma}
\label{lemma-flat-module-powers}
Let $R \to S$ be a ring map. Let $I \subset R$ be an ideal.
Let $M$ be an $S$-module. Assume
\begin{enumerate}
\item $R$ is a Noetherian ring,
\item $S$ is a Noetherian ring,
\item $M$ is a finite $S$-module, and
\item for each $n \geq 1$ the module $M/I^n M$ is flat over
$R/I^n$.
\end{enumerate}
Then for every $\mathfrak q \in V(IS)$
the localization $M_{\mathfrak q}$ is flat over $R$.
In particular, if $S$ is local and $IS$ is contained
in its maximal ideal, then $M$ is flat over $R$.
\end{lemma}

\begin{proof}
We are going to use
Lemma \ref{lemma-variant-local-criterion-flatness}.
By assumption $M/IM$ is flat over $R/I$. Hence it suffices to check
that $\text{Tor}_1^R(M, R/I)$ is zero on localization at $\mathfrak q$. By
Remark \ref{remark-Tor-ring-mod-ideal}
this Tor group is equal to $K = \Ker(I \otimes_R M \to M)$.
We know that the kernel of
$I/I^n \otimes_{R/I^n} M/I^nM \to M/I^nM$ is zero for all $n \geq 1$.
Hence an element of $K$ maps to zero in $I/I^n \otimes_{R/I^n} M/I^nM$.
Since
$$
I/I^n \otimes_{R/I^n} M/I^nM = I/I^n \otimes_R M =
(I \otimes_R M)/I^{n - 1}(I \otimes_R M)
$$
we conclude that $K \subset I^{n - 1}(I \otimes_R M)$ for all $n \geq 1$.
By the Artin-Rees lemma, and more precisely
Lemma \ref{lemma-intersection-powers-ideal-module}
we conclude that $K_{\mathfrak q} = 0$, as desired.
\end{proof}

\begin{lemma}
\label{lemma-surjective-on-tor-one}
Let $R \to R' \to R''$ be ring maps.
Let $M$ be an $R$-module. Suppose that $M \otimes_R R'$
is flat over $R'$. Then the natural map
$\text{Tor}_1^R(M, R') \otimes_{R'} R'' \to
\text{Tor}_1^R(M, R'')$ is onto.
\end{lemma}

\begin{proof}
Let $F_\bullet$ be a free resolution of $M$ over $R$.
The complex $F_2 \otimes_R R' \to F_1\otimes_R R' \to F_0 \otimes_R R'$
computes $\text{Tor}_1^R(M, R')$.
The complex $F_2 \otimes_R R'' \to F_1\otimes_R R'' \to F_0 \otimes_R R''$
computes $\text{Tor}_1^R(M, R'')$. Note that
$F_i \otimes_R R' \otimes_{R'} R'' = F_i \otimes_R R''$. Let
$K' = \Ker(F_1\otimes_R R' \to F_0 \otimes_R R')$ and
similarly $K'' = \Ker(F_1\otimes_R R'' \to F_0 \otimes_R R'')$.
Thus we have an exact sequence
$$
0 \to K' \to F_1\otimes_R R' \to F_0 \otimes_R R' \to M \otimes_R R' \to 0.
$$
By the assumption that $M \otimes_R R'$ is flat over $R'$,
the sequence
$$
K' \otimes_{R'} R'' \to
F_1 \otimes_R R'' \to
F_0 \otimes_R R'' \to
M \otimes_R R'' \to 0
$$
is still exact. This means that $K' \otimes_{R'} R'' \to K''$
is surjective. Since $\text{Tor}_1^R(M, R')$ is a quotient of $K'$ and
$\text{Tor}_1^R(M, R'')$ is a quotient of $K''$ we win.
\end{proof}

\begin{lemma}
\label{lemma-surjective-on-tor-one-trivial}
Let $R \to R'$ be a ring map. Let $I \subset R$ be
an ideal and $I' = IR'$. Let $M$ be an $R$-module
and set $M' = M \otimes_R R'$. The natural map
$\text{Tor}_1^R(R'/I', M) \to \text{Tor}_1^{R'}(R'/I', M')$
is surjective.
\end{lemma}

\begin{proof}
Let $F_2 \to F_1 \to F_0 \to M \to 0$ be a free resolution of
$M$ over $R$. Set $F_i' = F_i \otimes_R R'$. The sequence
$F_2' \to F_1' \to F_0' \to M' \to 0$ may no longer be exact
at $F_1'$. A free resolution of $M'$ over $R'$ therefore looks
like
$$
F_2' \oplus F_2'' \to F_1' \to F_0' \to M' \to 0
$$
for a suitable free module $F_2''$ over $R'$. Next, note that
$F_i \otimes_R R'/I' = F_i' / IF_i' = F_i'/I'F_i'$.
So the complex $F_2'/I'F_2' \to F_1'/I'F_1' \to F_0'/I'F_0'$
computes $\text{Tor}_1^R(M, R'/I')$. On the other hand
$F_i' \otimes_{R'} R'/I' = F_i'/I'F_i'$ and similarly
for $F_2''$. Thus the complex
$F_2'/I'F_2' \oplus F_2''/I'F_2'' \to F_1'/I'F_1' \to F_0'/I'F_0'$
computes $\text{Tor}_1^{R'}(M', R'/I')$. Since the vertical
map on complexes
$$
\xymatrix{
F_2'/I'F_2' \ar[r] \ar[d] &
F_1'/I'F_1' \ar[r] \ar[d] &
F_0'/I'F_0' \ar[d] \\
F_2'/I'F_2' \oplus F_2''/I'F_2'' \ar[r] &
F_1'/I'F_1' \ar[r] &
F_0'/I'F_0'
}
$$
clearly induces a surjection on cohomology we win.
\end{proof}

\begin{lemma}
\label{lemma-another-variant-local-criterion-flatness}
Let
$$
\xymatrix{
S \ar[r] & S' \\
R \ar[r] \ar[u] & R' \ar[u]
}
$$
be a commutative diagram of local homomorphisms of local Noetherian rings.
Let $I \subset R$ be a proper ideal.
Let $M$ be a finite $S$-module.
Denote $I' = IR'$ and $M' = M \otimes_S S'$.
Assume that
\begin{enumerate}
\item $S'$ is a localization of the tensor product
$S \otimes_R R'$,
\item $M/IM$ is flat over $R/I$,
\item $\text{Tor}_1^R(M, R/I) \to \text{Tor}_1^{R'}(M', R'/I')$
is zero.
\end{enumerate}
Then $M'$ is flat over $R'$.
\end{lemma}

\begin{proof}
Since $S'$ is a localization of $S \otimes_R R'$ we see that
$M'$ is a localization of $M \otimes_R R'$. Note that
by Lemma \ref{lemma-flat-base-change} the module $M/IM \otimes_{R/I} R'/I'
= M \otimes_R R' /I'(M \otimes_R R')$ is flat over $R'/I'$. Hence also
$M'/I'M'$ is flat over $R'/I'$ as the localization of a flat module
is flat. By Lemma \ref{lemma-variant-local-criterion-flatness}
it suffices to show that $\text{Tor}_1^{R'}(M', R'/I')$ is zero.
Since $M'$ is a localization of $M \otimes_R R'$, the last assumption
implies that it suffices to show that
$\text{Tor}_1^R(M, R/I) \otimes_R R'
\to
\text{Tor}_1^{R'}(M \otimes_R R', R'/I')$
is surjective.

\medskip\noindent
By Lemma \ref{lemma-surjective-on-tor-one-trivial} we see that
$\text{Tor}_1^R(M, R'/I') \to \text{Tor}_1^{R'}(M \otimes_R R', R'/I')$
is surjective. So now it suffices to show that
$\text{Tor}_1^R(M, R/I) \otimes_R R'
\to
\text{Tor}_1^R(M, R'/I')$
is surjective. This follows from Lemma \ref{lemma-surjective-on-tor-one}
by looking at the ring maps $R \to R/I \to R'/I'$ and the module $M$.
\end{proof}

\noindent
Please compare the lemma below to
Lemma \ref{lemma-criterion-flatness-fibre-nilpotent}
(the case of a nilpotent ideal) and
Lemma \ref{lemma-criterion-flatness-fibre}
(the case of finitely presented algebras).

\begin{lemma}[Crit\`ere de platitude par fibres; Noetherian case]
\label{lemma-criterion-flatness-fibre-Noetherian}
Let $R$, $S$, $S'$ be Noetherian local rings and let $R \to S \to S'$
be local ring homomorphisms. Let $\mathfrak m \subset R$ be the
maximal ideal. Let $M$ be an $S'$-module. Assume
\begin{enumerate}
\item The module $M$ is finite over $S'$.
\item The module $M$ is not zero.
\item The module $M/\mathfrak m M$
is a flat $S/\mathfrak m S$-module.
\item The module $M$ is a flat $R$-module.
\end{enumerate}
Then $S$ is flat over $R$ and $M$ is a flat $S$-module.
\end{lemma}

\begin{proof}
Set $I = \mathfrak mS \subset S$. Then we see that $M/IM$ is a flat
$S/I$-module because of (3). Since
$\mathfrak m \otimes_R S' \to I \otimes_S S'$ is surjective we see
that also $\mathfrak m \otimes_R M \to I \otimes_S M$ is surjective.
Consider
$$
\mathfrak m \otimes_R M \to I \otimes_S M \to M.
$$
As $M$ is flat over $R$ the composition is injective
and so both arrows are injective.
In particular $\text{Tor}_1^S(S/I, M) = 0$ see
Remark \ref{remark-Tor-ring-mod-ideal}. By
Lemma \ref{lemma-variant-local-criterion-flatness} we conclude
that $M$ is flat over $S$. Note that since $M/\mathfrak m_{S'}M$
is not zero by Nakayama's Lemma \ref{lemma-NAK}
we see that actually $M$ is faithfully flat over $S$ by
Lemma \ref{lemma-ff} (since it forces $M/\mathfrak m_SM \not = 0$).

\medskip\noindent
Consider the exact sequence
$0 \to \mathfrak m \to R \to \kappa \to 0$.
This gives an exact sequence
$0 \to \text{Tor}_1^R(\kappa, S) \to \mathfrak m \otimes_R S \to I \to 0$.
Since $M$ is flat over $S$ this gives an exact sequence
$0 \to \text{Tor}_1^R(\kappa, S)\otimes_S M \to
\mathfrak m \otimes_R M \to I \otimes_S M \to 0$.
By the above this implies that $\text{Tor}_1^R(\kappa, S)\otimes_S M = 0$.
Since $M$ is faithfully flat over $S$ this implies that
$\text{Tor}_1^R(\kappa, S) = 0$ and we conclude that
$S$ is flat over $R$ by Lemma \ref{lemma-local-criterion-flatness}.
\end{proof}

\begin{lemma}
\label{lemma-flatness-scallop}
Let $A$ be a ring, let $M$ be an $A$-module, and let $f \in A$. If
\begin{enumerate}
\item $f$ is a nonzerodivisor on $A$ and $M$,
\item $M_f$ is a flat $A_f$-module, and
\item $M/fM$ is a flat $A/fA$-module,
\end{enumerate}
Then $M$ is a flat $A$-module. Same with ``flat'' replaced by
``faithfully flat''.
\end{lemma}

\begin{proof}
Since $0 \to A \to A \to A/fA \to 0$ and $0 \to M \to M \to M/fM \to 0$ are
exact, we find that $\text{Tor}_i^A(M, A/fA) = 0$ for $i = 1$ and $i = 2$.
By Lemma \ref{lemma-what-does-it-mean} we conclude that
$\text{Tor}_1^A(M, N) = 0$ for all $A$-modules $N$ annihilated by $f$
(this uses the flatness of $M/fM$ over $A/fA$).
Given an $A$-module $N$ annihilated by $f$ we may
choose a short exact sequence $0 \to N' \to F \to N \to 0$ of $A$-modules
where $F$ is a direct sum of copies of $A/fA$. From the exact sequence
$$
\text{Tor}_2^A(M, F) \to
\text{Tor}_2^A(M, N) \to
\text{Tor}_1^A(M, N')
$$
we conclude that $\text{Tor}_2^A(M, N) = 0$ for all $A$-modules
$N$ annihilated by $f$. Next, let $K$ be an arbitrary $A$-module.
We may break the map $f : K \to K$ into two short exact sequences
$$
0 \to K[f] \to K \to K' \to 0
\quad\text{and}\quad
0 \to K' \to K \to K/fK \to 0
$$
where $K' = K/K[f] \cong fK$. Applying the exact sequences of Tor we
obtain exact sequences
$$
\text{Tor}_1^A(K[f], M) \to
\text{Tor}_1^A(K, M) \to
\text{Tor}_1^A(K', M)
$$
and
$$
\text{Tor}_2^A(K/fK, M) \to
\text{Tor}_1^A(K', M) \to
\text{Tor}_1^A(K, M)
$$
Using the vanishing of
$\text{Tor}_1^A(K[f], M)$ and
$\text{Tor}_2^A(K/fK, M)$
we conclude that $f : K \to K$ induces an injective map on
$\text{Tor}^A_1(M, K)$. In other words, we see that $\text{Tor}^A_1(M, K)$
is zero if and only if $\text{Tor}^A_1(M, K) \otimes_A A_f$ is zero.
By Lemma \ref{lemma-flat-base-change-tor} we have
$$
\text{Tor}^A_1(M, K) \otimes_A A_f = \text{Tor}^{A_f}_1(M_f, K_f)  = 0
$$
The vanishing by the flatness of $M_f$ over $A_f$
(Lemma \ref{lemma-characterize-flat}).
We conclude that $\text{Tor}_1^A(M, K) = 0$
for all $A$-modules $K$. Hence
$M$ is flat over $A$ by Lemma \ref{lemma-characterize-flat}.

\medskip\noindent
We omit the argument for the case of faithfully flat modules.
\end{proof}

\begin{lemma}
\label{lemma-flatness-scallop-pre}
Let $A$ be a ring, let $M$ be an $A$-module. Let $I = (f_1, \ldots, f_r)$
be an ideal of $A$ generated by $r \geq 1$ elements. If
\begin{enumerate}
\item $M_{f_i}$ is a flat $A_{f_i}$-module for $i = 1, \ldots, r$,
\item $M/IM$ is a flat $A/I$-module,
\item $\text{Tor}_i^A(M, A/I) = 0$ for $i = 1, \ldots, r + 1$.
\end{enumerate}
Then $M$ is a flat $A$-module. Same with ``flat'' replaced by
``faithfully flat''.
\end{lemma}

\begin{proof}
From Lemma \ref{lemma-what-does-it-mean} we see that
$\text{Tor}_1^A(M, K) = 0$ for all $A$-modules $K$ annihilated by $I$.
Given an $A$-module $K$ annihilated by $I$ we may choose a short
exact sequence $0 \to N \to F \to K \to 0$ of $A$-modules where $F$
is a direct sum of copies of $A/I$. We obtain an exact sequence
$$
\text{Tor}_i^A(M, F) \to
\text{Tor}_i^A(M, K) \to
\text{Tor}_{i - 1}^A(M, N)
$$
Thus using assumption (3) and induction on $i$
we conclude that $\text{Tor}_i^A(M, K) = 0$
for all $A$-modules $K$ annihilated by $I$ and $i = 1, \ldots, r + 1$.

\medskip\noindent
Suppose that for some $1 \leq j \leq r$
we have shown that $\text{Tor}_i^A(M, K) = 0$
for all $A$-modules $K$ annihilated by $f_1, \ldots, f_j$
and $i = 1, \ldots, j + 1$.
Let $K$ be an $A$-module annihilated by $f_1, \ldots, f_{j - 1}$.
We may break the map $f_j : K \to K$ into two short exact sequences
$$
0 \to K[f_j] \to K \to K' \to 0
\quad\text{and}\quad
0 \to K' \to K \to K/f_jK \to 0
$$
where $K' = K/K[f_j] \cong f_jK$. Let $1 \leq i \leq j$.  Applying the exact
sequences of Tor we obtain exact sequences
$$
\text{Tor}_i^A(K[f_j], M) \to
\text{Tor}_i^A(K, M) \to
\text{Tor}_i^A(K', M)
$$
and
$$
\text{Tor}_{i + 1}^A(K/f_jK, M) \to
\text{Tor}_i^A(K', M) \to
\text{Tor}_i^A(K, M)
$$
Using the vanishing of $\text{Tor}_i^A(K[f_j], M)$ and
$\text{Tor}_{i + 1}^A(K/f_jK, M)$ we conclude that $f_j : K \to K$
induces an injective map on $\text{Tor}^A_i(M, K)$.
In other words, we see that $\text{Tor}^A_i(M, K)$
is zero if and only if $\text{Tor}^A_i(M, K) \otimes_A A_{f_j}$ is zero.
By Lemma \ref{lemma-flat-base-change-tor} we have
$$
\text{Tor}^A_i(M, K) \otimes_A A_{f_j} =
\text{Tor}^{A_{f_j}}_i(M_{f_j}, K_{f_j})  = 0
$$
The vanishing by the flatness of $M_{f_j}$ over $A_{f_j}$
(Lemma \ref{lemma-characterize-flat}).
We conclude that $\text{Tor}_i^A(M, K) = 0$
for all $A$-modules $K$ annihilated by $f_1, \ldots, f_{j - 1}$
and $i = 1, \ldots, j$.
By descending induction on $j$, we conclude that
this holds for $j = 0$, i.e., we see that
$\text{Tor}_1^A(M, K) = 0$
for all $A$-modules $K$ . Hence
$M$ is flat over $A$ by Lemma \ref{lemma-characterize-flat}.

\medskip\noindent
We omit the argument for the case of faithfully flat modules.
\end{proof}








